% !PRINTED 599
\weberpage{599}
%<<CONT
\weq{6}{\alpha_0 = \alpha + M\vartheta \equiv 0 \pmod{\mathfrak{m}}}
bestimmen, worin $\vartheta$ eine ganze Zahl in $\Omega$ ist.

Daraus folgt, daß jedes im Exkludenten aufgehende Primideal
 $\mathfrak{e}$ in $M$ aufgehen muß, weil sonst $\alpha_0$ für $\mathfrak{m} = \mathfrak{e}$ durch $\mathfrak{e}$
teilbar wäre.

Haben $\alpha$ und $M$ einen gemeinschaftlichen Teiler, so haben
alle Zahlen der Form $\alpha + M\xi$, in der $\xi$ eine beliebige Zahl ist,
denselben gemeinschaftlichen Teiler mit $M$. Es gibt aber in $A$
auch Zahlen, die relativ prim zu $M$ sind, z.\,B. die Zahl 1. Es
gilt dann der folgende Satz:

3. \gesperrt{Ist $\mathfrak{m}$ relativ prim zu $M$, so kann man in $A$
eine durch $\mathfrak{m}$ teilbare ganze Zahl $\alpha_0 = \mathfrak{m}\mathfrak{n}$ von der
Art bestimmen, daß $\mathfrak{n}$ zu $M$ und einem beliebig
gegebenen Ideal $\mathfrak{c}$ relativ prim wird.}

Hat man nämlich nach (6) eine durch $\mathfrak{m}$ teilbare Zahl
\[
\alpha_0 = \mathfrak{m}\mathfrak{n}
\]
bestimmt, die zu $M$ teilerfremd ist, so erhält man andere solche
Zahlen:
\[
\alpha_0' = \alpha_0 + M\xi = \mathfrak{m}\mathfrak{n}',
\]
wenn
\[
\xi = \mathfrak{m}\mathfrak{r}
\]
eine beliebige durch $\mathfrak{m}$ teilbare ganze Zahl ist. Ist nun $\mathfrak{q}$ ein in $\mathfrak{c}$,
aber nicht in $\alpha_0$ aufgehendes Primideal, so nehme man $\mathfrak{r}$ durch
$\mathfrak{q}$ teilbar an, dann ist $\mathfrak{n}'$ nicht durch $\mathfrak{q}$ teilbar. Geht aber $\mathfrak{q}$ in $\mathfrak{c}$
und in $\alpha_0$ auf, so nehme man $\mathfrak{r}$ durch $\mathfrak{q}$ nicht teilbar an; dann
ist auch $\mathfrak{n}'$ nicht durch $\mathfrak{q}$ teilbar.

Daraus folgt:

4. \gesperrt{Ist $\mathfrak{m}$ ein beliebiges zu $M$ teilerfremdes Ideal,
so gibt es in jeder Idealklasse $\overline{A_\nu}$ Ideale, die zu $\mathfrak{m}$
relativ prim sind.}

Man nehme, um dies zu beweisen, ein Ideal $\mathfrak{b}$ in der Klasse
$(\overline{A_\nu})^{-1}$ und bestimme nach 3. eine Zahl
\[
\alpha = \mathfrak{a}\mathfrak{b},
\]
so daß $\mathfrak{a}$ relativ prim zu $\mathfrak{m}$ ist, dann gehört $\mathfrak{a}$ in die Klasse $\overline{A_\nu}$
und genügt der Forderung des Satzes 4.

\noindent\textbf{\S\ 162.\qquad \gesperrt{Die durch ein Ideal teilbaren Ideale der Hauptklasse.}}

\gesperrt{Es soll jetzt unter $\Omega$ der quadratische Körper mit
negativer Grundzahl $\Delta$ verstanden werden.}
